\scanpage{101}

\textcircled{\scriptsize $3^\circ$} On a déjà prouvé (14), il faut prouver (15) (16).

\uline{Formule} (15) : Pour $j<i$ cela provient de $\tau_j=\mathrm{id}+a'_j\otimes a_j$, et $a_j\in\unsure{E_i|\mathcal E_i}$ \inserted{(vrai sous les hyp. A) B))}, pour $j>i$ cela provient de $a'_j\mid E_i=0$, qui implique $E_i\subset\underset{\substack{\Vert\\(E^{\tau_j})}}{H_j}$.

\uline{Formule} 16\qquad On a

\[E_i=E_{i-1}+ka_{i-1}\]

\[\begin{aligned}\text{donc}\qquad\tau_i E_i&=\underset{\Vert}{\uline{\tau_i E_{i-1}}}+k\tau_i a_{i-1}\\ &=E_{i-1}+k\underset{\substack{\rotatebox{-90}{$\in$}\\ k^*}}{\underline{\langle a_{i-1},a'_i\rangle}}a_i=E_{i-1}+k\deleted{\unsure{E}}a_i\end{aligned}\]

et comme $(a_{i-1},\deleted{\unsure{et}}a_i)$ \deleted{\unsure{sont}} est libre mod $E_{i-1}$, cela prouve notre assertion.

\uline{Corollaire 2} Sous les conditions du \deleted{\unsure{th.}} \inserted{corollaire} \inserted{(a fortiori sous les conditions du th 1 2°)}) \deleted{\unsure{on aura|on a}}

\[\deleted{\unsure{H_i\overset{\mathrm{dfn}}=|H_i\overset{\mathrm{df}}=}}\quad\left\{\begin{aligned}E^{\tau_i}&=\operatorname{Ker}a'_i\ (\overset{\mathrm{dfn}}=H_i)\\ka_i&=\operatorname{Im}(\tau_i-\mathrm{id})\ (\overset{\mathrm{dfn}}=h_i)\end{aligned}\right.\quad\text{pour}\quad0\leq i\leq n\tag*{(21)}\]

\[E_i\subset H_j\quad\deleted{\unsure{(\ =E^{\tau_j})\ \operatorname{Ker}a'_i|(\operatorname{Ker}a'_j)}}\qquad\text{pour }0\leq i<j\leq n\deleted{\unsure{a}}\tag*{(22)}\]

\[E_i\cap H_i=E_{i-1},\tag*{(23)}\]

\[\tau_i\tau_{i+1}\ne\tau_{i+1}\tau_i,\quad\uline{\text{même modulo }k^*}\qquad\deleted{\unsure{\text{si}}}\ (\text{si }k\ne0)\tag*{(24)}\]

\[\deleted{\underset{\substack{\Vert\\ka_i}}{h_i}\subset\underset{\substack{\Vert\\\operatorname{Ker}a'_{i+1}}}{H_{i+1}}}\quad\text{i.e.}\quad\langle a_i,a'_{i+1}\rangle\in k^*\ \inserted{-1\leq i\leq n}\ (\text{donc }\ne0\text{ si }k\ne0)\tag*{(25)}\]

\[\underset{\substack{\Vert\\ka_i}}{h_i}\not\subset\underset{\substack{\Vert\\\operatorname{Ker}a'_{i+1}}}{H_{i+1}}\qquad\text{si}\quad k\ne0\tag*{(25$'$)}\]

La formule \unsure{(22)|(21)} est déjà prouvée, (21) 1° provient du fait que $\{a_i\}$ est libre, (21 2°) du fait que $a'_i:\deleted{\unsure{\text{est surj}}} E\longrightarrow k$ est surjectif, \deleted{\unsure{ceci}} car $\langle a_{i-1},a'_i\rangle\in k^*$ (NB $a_{-1}=s_0$ !) grâce à (13) (4). \deleted{On a} La formule (23) provient de $E_i=k(a_{-1},\ldots,a_{i-1})$, et que $H_i=\operatorname{Ker}a'_i$, $a'_i$ \unsure{nul|nulle} sur $a_j$ $(-1\leq j\leq i-2)$, $\langle a_{i-1},a'_i\rangle\in k^*$.

% archive sheet number: 100
